% Part: proof-theory
% Chapter: cut-elimination
% Section: introduction

\documentclass[../../../include/open-logic-section]{subfiles}

\begin{document}

\olfileid{pt}{cut}{int}

\olsection{परिचय}

\Cut{} नियम असा आहे:
\[
\Axiom$\Gamma \fCenter \Delta, !A$
\Axiom$!A, \Pi \fCenter \Lambda$
\RightLabel{\Cut}
\BinaryInf$\Gamma, \Pi \fCenter \Delta, \Lambda$
\DisplayProof
\]
(!!{formula}~$!A$ याला
\emph{कट-सूत्र} म्हणतात.)

\Cut{} हा नियम गेंटझेन यांच्या
मूळ क्रमवर्ती कलनात~\Log{LK} होता.
गेंटझेन यांचे \emph{Hauptsatz}
(“मुख्य प्रमेय”) असे सांगते की
~\Log{LK} मध्ये !!{provable} असलेली
प्रत्येक क्रमवर्ती \Cut{} न वापरताही
!!{provable} आहे; म्हणजे
~\Log{LK} मधील सिद्धतांतून
\Cut{} “काढून टाकता” येतो.
आपल्या \Log{G1c} आणि \Log{G3c}
पद्धतींत \Cut{} नाही; तरी त्यांच्याबाबतही
प्रश्न पडतो: \Cut{} आणि
~\Log{G1c} चे (किंवा \Log{G3c} चे)
इतर नियम वापरणाऱ्या !!{proof}s मधून
\Cut{} काढून टाकता येईल का?
आपल्या आधी निश्चित केलेल्या
परिभाषेत हाच प्रश्न असा आहे:
\Cut{} हा \Log{G1c} चा
(किंवा \Log{G3c} चा)
ग्राह्य नियम आहे का?

कट-निर्मूलन प्रमेय हा सिद्धता
उपपत्तीतील कदाचित सर्वाधिक महत्त्वाचा
आणि प्रसिद्ध निष्कर्ष आहे.
सुरुवातीला आवश्यक वाटू शकणारा
नियम \Log{G1c} आणि \Log{G3c}
यांना प्रत्यक्षात लागत नाही, हे
तो दाखवतो. कारण प्रत्यक्ष सिद्धता
औपचारिक करताना उपप्रमेयांचा
वापर कसा दाखवायचा हा प्रश्न
पडतो. गणितज्ञ प्रथम एक निष्कर्ष
सिद्ध करतात आणि मग दुसऱ्या
निष्कर्षाच्या सिद्धतेत त्याचा
आधार घेतात. म्हणून उपप्रमेय
$L$ ची सिद्धता प्रथम
$\Sequent L$ या क्रमवर्तीच्या
!!{proof} म्हणून औपचारिक करता येईल.
मग $L$ वापरणाऱ्या दुसऱ्या
निष्कर्ष~$R$ ची सिद्धता
$L \Sequent
R$ या क्रमवर्तीच्या
!!{proof} म्हणून औपचारिक करू.
केवळ $R$ ची, म्हणजे
$\Sequent R$ ची
!!a{proof} आहे, हे कसे ठरवायचे?
त्या दोन !!{proof}s जोडण्याचा
मार्ग हवा; \Cut{} नियमाने
ते करता येईल.

\Log{G1c} आणि \Log{G3c}
संपूर्ण आहेत, असे आपण सांगितले:
म्हणून सिद्ध होण्याची अपेक्षा
असलेल्या सर्व वैध क्रमवर्ती
ते सिद्ध करतात. हे थेटही
सिद्ध करता येते. पण गेंटझेन
लिहीत होते तेव्हा केवळ स्वयंसिद्धकीय
निष्पत्ती पद्धती संपूर्ण आहेत,
हे ज्ञात होते. \Log{LK} ची
संपूर्णता दाखवण्यासाठी गेंटझेन
यांनी स्वयंसिद्धकीय पद्धतीतील
!!{derivation}s चे
~\Log{LK} मधील !!{proof}s मध्ये
भाषांतर दिले; त्यात
\Cut{} अत्यावश्यक होता.
कट-निर्मूलन केवळ अभिजात
तर्कशास्त्रापुरते महत्त्वाचे नाही:
इतर अनेक तर्कशास्त्रांनाही
क्रमवर्ती कलने आहेत. काहींसाठी
त्या कलनांची संपूर्णता थेट
सिद्ध करणे कठीण किंवा अशक्य
असते; त्यामुळे इतर पद्धतींमधून
\Cut{} वापरून केलेले भाषांतर
हाच संपूर्णता दाखवण्याचा मार्ग असतो.
मग \Cut{} अनावश्यक आहे आणि
\Cut{} नसलेली पद्धतही संपूर्ण
आहे, हे स्थापित करण्यासाठी
सिद्धता-उपपत्तीय कट-निर्मूलन
सिद्ध करणे आवश्यक ठरते.

कट-निर्मूलनाचा उपयोग इतर
स्वतंत्रपणे महत्त्वाचे निष्कर्ष
मिळवण्यासाठीही होतो. त्यांपैकी
क्रेगचे अंतर्वेशन उपप्रमेय आणि
एरब्राँचे प्रमेय ही दोन उदाहरणे
आपण पाहू.

कट-निर्मूलनाच्या सिद्धतेत
बहुधा \Cut{} वापरणाऱ्या
!!a{proof} चे त्याविना
सिद्धतेत रूपांतर कसे करावे,
हे दाखवले जाते. ही रूपांतरे
सहसा “स्थानिक” असतात:
सिद्धतेचा एक भाग
(साधारणपणे \Cut{} अनुमानाने
संपणारी उप-!!{proof})
घेऊन त्याच्या जागी
त्याच क्रमवर्तीची दुसरी
उप-!!{proof} ठेवतो;
तिच्यातील \Cut{} अनुमान
काढलेले किंवा इतर
अनुमानांनी बदललेले असते.
अशा युक्तिवादांनी \Cut{}
असलेल्या !!{proof}s चे
त्याविना सिद्धतांत रूपांतर
करणारी पायरी-पायरीची
अल्गोरिदम मिळते. उपयोगांत
कट-निर्मूलन प्रक्रिया वापरताना
यातून अधिक माहिती मिळते.
उदाहरणार्थ, अंतर्वेशन उपप्रमेयाची
अशी प्रतिमान-उपपत्तीय सिद्धता
आहे की “$!A \lif !B$”
वैध असेल, तर अंतर्वेशक अस्तित्वात
असतो (आत्ता अंतर्वेशक म्हणजे
काय याकडे दुर्लक्ष करा).
कट-निर्मूलन वापरणारा
सिद्धता-उपपत्तीय युक्तिवाद
$!A
\lif !B$ ची !!a{proof}
घेऊन अंतर्वेशक देणारी
अल्गोरिदम पुरवतो. प्रतिमान
उपपत्तीच्या युक्तिवादाच्या
तुलनेत ही सिद्धता-उपपत्तीय
पद्धत \emph{रचनात्मक} आहे.

कट-निर्मूलन प्रमेय सिद्ध
करण्याच्या दोन प्रचलित पद्धती
आहेत. पहिली, गेंटझेनपासून
आलेली, सिद्धतेतील सर्वोच्च
स्थानी असलेले कट काढता
येतात हे दाखवते आणि मग
असे कट संपेपर्यंत काढत जाते.
दुसरी, मूळची श्यूटे आणि
टेट यांची, !!a{proof} मधील
सर्वाधिक गुंतागुंतीच्या कटांवर
लक्ष केंद्रित करते; इतर
कोणत्याही कटच्या !!{formula}
ची खोली अधिक नाही, या
अर्थाने महत्तम असलेले
कट ती एकामागोमाग काढते.
दोन्ही पद्धतींचे फायदे
आणि मर्यादा आहेत. काही
पद्धतींमध्ये एक मार्ग
चालतो, तर दुसरा अवघड
किंवा अशक्य असतो.
म्हणून आपण दोन्ही मार्ग
वर्णन करू. ~\Log{G3c}
साठी कट-निर्मूलन सिद्ध करू.
खरे तर नेहमीचा \Cut{}
नव्हे, तर \emph{संदर्भ-सामायिक कट}~\CutCS:
\[
\Axiom$\Gamma \fCenter \Delta, !A$
\Axiom$!A, \Gamma \fCenter \Delta$
\RightLabel{\Cut}
\BinaryInf$\Gamma \fCenter \Delta$
\DisplayProof
\]
एखाद्या क्रमवर्ती कलनात
दुर्बलीकरण आणि आकुंचन
हे प्रत्यक्ष नियम म्हणून
(जसे \Log{G1c} मध्ये)
किंवा ग्राह्य नियम म्हणून
(जसे \Log{G3c} मध्ये)
असतील, तर \Cut{} द्वारे
\CutCS चे अनुकरण करता
येते, आणि उलटही.
आपण \Cut ऐवजी \CutCS{}
निवडतो, कारण पहिली
कट-निर्मूलन पद्धत \CutCS साठी
वर्णन करणे सोपे आहे,
आणि दुसरी पद्धत केवळ
\CutCS साठीच चालते.

\end{document}
